\documentclass[12pt]{article}
\usepackage[utf8]{inputenc}
\usepackage[T1]{fontenc}
\usepackage[margin=1in]{geometry}
\usepackage{setspace}
\usepackage{xcolor}
\usepackage[colorlinks=true, urlcolor=blue!60!black, linkcolor=blue!60!black]{hyperref}
\setlength{\parindent}{0em}
\setlength{\parskip}{0.6em}
\renewcommand{\baselinestretch}{1.0}
\begin{document}
\begin{center}
{\LARGE \textit{Champion v. Ames}, 188 U.S. 321 (1903)\par}
\vspace{0.5em}
{\large Case Brief --- Two-Page Edition\par}
\end{center}

\section{Memory Jogger}
The Lottery Case: lottery tickets are articles of commerce, and the power to regulate interstate commerce includes the power to prohibit it (\href{https://supreme.justia.com/cases/federal/us/188/321/}{Justia}).

\section{Detailed Case Facts}
The Federal Lottery Act of 1895 made it a crime to cause any lottery ticket to be ``carried from one state to another.'' Charles F. Champion was indicted for conspiring to ship two Paraguayan lottery tickets by Wells-Fargo Express from Dallas to Fresno, California.

\section{Procedural History}
Held by U.S. Marshal John C. Ames, Champion sought habeas corpus, arguing that the Act exceeded the commerce power. The circuit court dismissed the writ. After three arguments, the Supreme Court affirmed in 1903.

\section{Judicial Votes}
5--4. Harlan wrote for Brown, White, McKenna, and Holmes; Chief Justice Fuller dissented, joined by Brewer, Shiras, and Peckham.

\section{Holding}
Lottery tickets are subjects of traffic and of commerce, and Congress may prohibit their interstate carriage as injurious to public morals without violating the Tenth or Fifth Amendment.

\section{Analysis of Opinions}
\textbf{Majority:} Under \textit{Gibbons}, commerce includes ``traffic'' and ``the transit of persons''; the tickets had a market price and were carried for hire. Under \textit{McCulloch}, Congress may choose prohibition, as with diseased cattle and liquor. A delegated power cannot invade reserved powers, and Congress merely supplemented state lottery bans. \textbf{Fuller:} suppressing lotteries belongs to the States' police power; tickets are no more commerce than ``an invitation to dine,'' and they ``cannot communicate bad principles by contact.'' \textbf{Later:} \textit{Darby} (1941) overruled \textit{Hammer v. Dagenhart} and vindicated \textit{Champion}.

\section{Examples: Future Applications}
\textbf{Same-Side:} Congress may make it a crime to ship unserialized, build-it-yourself firearm kits across state lines. \textbf{Opposite-Side:} a federal ban on public intoxication with no interstate element fails under \textit{Lopez}. \textbf{Fence-Sitter:} punishing possession of a device that once crossed state lines.

\section{Critique}
As Fuller foresaw, once Congress may close channels for moral reasons, the ``harmful-in-itself'' limit collapses. Banning interstate shipment also lets Congress reach local production in fact while disclaiming a police power.

\section{Key Quotations}
\begin{quote}
``As a state may, for the purpose of guarding the morals of its own people, forbid all sales of lottery tickets within its limits, so Congress, for the purpose of guarding the people of the United States against the `widespread pestilence of lotteries' and to protect the commerce which concerns all the states, may prohibit the carrying of lottery tickets from one state to another.''
\end{quote}
\noindent --- Justice Harlan, for the Court, 188 U.S. at 357.

\end{document}
